\section{Evidence-Allocation Framework and Legal Scope}
\label{app:framework_submission}

The model used in the main text is intentionally spare. Its job
is to justify an ordering rule for scarce forensic attention,
not to model the full internal organization of a cartel. Issues
such as side payments, discipline, and leniency incentives are
treated in the cartel literature
\citep{asker2010leniency,marshall2012economics}. The narrower
object here is the pre-proof enforcement problem: why a firm
that repeatedly loses, while continuing to appear in procurement
records, can be a useful target for bid-level follow-up.

\subsection{Primitives}
\label{app:framework_primitives}

Consider a procurement environment with repeated tender-items.
In a cartel-targeted tender, one firm is assigned the winning
role and submits the designated bid $b^*$. A losing-role
participant submits a bid $b_\ell>b^*$ and is not meant to win.
The losing role has a cost $c_1>0$. Winning by mistake or by
deviation creates an internal cartel sanction $\kappa>0$,
interpreted broadly as the loss of cartel discipline, future
rents, or capacity-allocation rights. Let $r$ denote the private
gain from deviating into the winner role in a tender where the
firm was assigned to lose.

Firms outside the winning role are of two reduced-form types.
Type $C$ firms are cartel-deployed losing-role participants. Type
$G$ firms are ordinary losing participants: they may be
inefficient, capacity constrained, inexperienced, or simply
unlucky. The type notation is a device for stating the maintained
ranking condition; it is not an empirical classification rule.
The empirical screen conditions on the observed event
$W_i=0$, where $W_i$ is the firm's number of wins, and ranks
firms by participation count $T_i$.

\begin{lemma}[Zero-win losing role]
\label{lem:submission_zero_win}
If the internal sanction from winning in a losing role exceeds
the private deviation gain, $\kappa>r$, a losing-role firm's
equilibrium bid lies above the designated winning bid:
$b_\ell>b^*$. Hence $\Pr(W_i>0\mid C)=0$ over tenders in which
the firm is deployed in the losing role.
\end{lemma}

\begin{proof}
If the losing-role participant submits $b_\ell>b^*$, it loses
and receives the deployment payoff net of the role cost. If it
deviates to a winning bid, it gains $r$ but pays the sanction
$\kappa$. The deviation is unprofitable when $\kappa>r$. The
result fixes the win probability of the losing role, not the
exact numerical bid. Any bid above the designated winner's bid is
consistent with the reduced-form role. 
\end{proof}

Lemma~\ref{lem:submission_zero_win} explains why the zero-win
conditioning set is a coherent place to begin. The conditioning
set is intentionally broad: many genuine firms also never win.
The evidentiary content comes from the second observable,
participation intensity.

\subsection{Participation as a Ranking Statistic}
\label{app:framework_participation}

Within the zero-win set, let participation counts follow a
Poisson approximation:
\[
  T_i \mid \tau_i, q \sim \text{Poisson}(\lambda_q \tau_i),
  \qquad q\in\{C,G\},
\]
where $\tau_i$ is the length of time firm $i$ is observed and
$\lambda_q$ is the type-specific participation rate. The
screening assumption is
\[
  \lambda_C>\lambda_G
  \quad\text{on the conditioning set } W_i=0.
\]
This is the substantive assumption behind the frequent-loser
ranking. Cartel-deployed losing-role participants, if present,
are repeatedly deployed to provide losing participation; ordinary
zero-win firms participate less often.
The condition $\lambda_C>\lambda_G$ is \emph{maintained}, not
derived: it is the behavioral assumption under which repeated
zero-win participation becomes a monotone ranking statistic for
loser-side exposure. The empirical sections then ask whether
the ranking implied by this assumption has validation content
against adjudication-adjacent labels.
The type notation is a device for stating the maintained ranking
condition; it is not an empirical classification rule.

\begin{proposition}[Monotone loser-side suspicion]
\label{prop:submission_monotone}
Under the Poisson approximation and $\lambda_C>\lambda_G$, the
posterior probability $\Pr(C\mid T_i=t,W_i=0,\tau_i)$ is
non-decreasing in $t$. Consequently, $T_i$ and
$\log(1+T_i)$ are rank-equivalent screening statistics for
loser-side cartel-adjacency exposure within the zero-win
stratum.
\end{proposition}

\begin{proof}
By Bayes' rule, the posterior odds of type $C$ against type $G$
are proportional to the likelihood ratio
\[
  \frac{\Pr(T_i=t\mid C,\tau_i)}
       {\Pr(T_i=t\mid G,\tau_i)}
  =
  \exp[-(\lambda_C-\lambda_G)\tau_i]
  \left(\frac{\lambda_C}{\lambda_G}\right)^t .
\]
Because $\lambda_C>\lambda_G$, the likelihood ratio is increasing
in $t$. This is the monotone-likelihood-ratio property of the
Poisson family \citep{karlin1956mlr}. The posterior increases
with participation count, and the log transformation
$\log(1+T_i)$ preserves the ranking.
\end{proof}

The empirical binary rule used in the main text is an operational
coarsening of Proposition~\ref{prop:submission_monotone}. It flags firms above the
median plus $1.5$ times the interquartile range of zero-win
participation. The threshold is not a structural parameter. It
is an auditable way to turn a monotone ranking into an
operational list. Agencies could instead use the continuous
score, a capacity-constrained top-$k$ list, or a threshold
calibrated to their investigative budget.
The posterior language belongs to the reduced-form screening
model. The empirical implementation uses the statistic as a
ranking device for forensic priority, not as an estimated
probability of legal membership.

\subsection{Evidentiary Scope}
\label{app:framework_scope}

The proposition delivers a ranking result. Even if the model is
correct, a high value of
$\log(1+T_i)$ raises the posterior probability of loser-side
cartel-adjacency exposure only relative to other zero-win firms.
It does not
identify an agreement, assign cartel membership, prove that any
specific bid was cover, or establish a damages base. Those
claims require evidence from the richer record: bid values,
communications, leniency material, repeated winner rotation, or
other case-specific facts.

This boundary is the legal reason for the staged architecture in
the main text. The award layer can allocate attention because it
ranks firms on a statistic that is cheap, observable, and
monotone under a transparent behavioral assumption. The bid layer
and case record then do the evidentiary work that liability
requires, because liability turns on conduct and agreement, not
on suspicious participation alone.

The framework does not imply that all zero-win firms occupy a
cartel losing role, that frequent-loser status is legal
membership, or that any specific bid is a cover bid. It supplies
only the ranking logic for triage.

\begin{proposition}[Scope separation]
\label{prop:submission_scope}
Under Lemma~\ref{lem:submission_zero_win} and
Proposition~\ref{prop:submission_monotone}, the frequent-loser statistic
is sufficient to order expected investigative value within the
zero-win stratum. It is not sufficient for liability or damages.
\end{proposition}

\begin{proof}
Triage requires an ordering of expected investigative value.
Proposition~\ref{prop:submission_monotone} supplies that ordering
within the zero-win stratum. Liability and damages require
different predicates: proof of agreement, conduct in particular
tenders, and an economic measure of price injury. The statistic contains
participation and win information only. It omits the facts needed
for those predicates. The same statistic can be useful for
allocating forensic attention and insufficient for legal
adjudication.
\end{proof}
